\section{The Hopf invariant}\label{ch:invariant}

This section proves Theorem~\ref{thm:B}. Section~\ref{sec:hdef} defines the invariant and proves the rules for composition that are used throughout; Section~\ref{sec:powers} computes the degree of the octonionic power maps.

\subsection{Definition and composition rules}\label{sec:hdef}

Let $n\ge2$ and $f\colon S^{2n-1}\to S^n$. The mapping cone $C_f=S^n\cup_fD^{2n}$ is a CW complex with cells in dimensions $0$, $n$ and $2n$, so
\[
H^n(C_f;\Z)\cong\Z\,\alpha_f,\qquad H^{2n}(C_f;\Z)\cong\Z\,\beta_f,
\]
where $\alpha_f$ restricts to the chosen generator of $H^n(S^n)$ and $\beta_f$ corresponds to the chosen generator of $H^{2n}(D^{2n},S^{2n-1})$ under the isomorphisms $H^{2n}(C_f)\cong H^{2n}(C_f,S^n)\cong H^{2n}(D^{2n},S^{2n-1})$.

\begin{definition}\label{def:hopfinv}
The \emph{Hopf invariant} $\Hopf(f)\in\Z$ is defined by $\alpha_f\smile\alpha_f=\Hopf(f)\,\beta_f$.
\end{definition}

If $f\simeq f'$, then $C_f$ and $C_{f'}$ are homotopy equivalent by a map that is the identity on $S^n$ and of degree one on the top cell, so $\Hopf(f)$ depends only on the homotopy class of $f$. Changing the orientation of $S^{2n-1}$ changes the sign of $\Hopf$, and changing that of $S^n$ does not. The function $\Hopf\colon\pi_{2n-1}(S^n)\to\Z$ is a homomorphism \cite[Proposition~4B.1]{Hat02}; we use this only in Corollary~\ref{cor:image} and in Section~\ref{ch:stable}.

\begin{proposition}\label{prop:hodd}
If $n$ is odd then $\Hopf(f)=0$ for every $f\colon S^{2n-1}\to S^n$.
\end{proposition}

\begin{proof}
By graded commutativity $\alpha_f^2=(-1)^{n^2}\alpha_f^2=-\alpha_f^2$, so $2\Hopf(f)\beta_f=0$ in $H^{2n}(C_f)\cong\Z$.
\end{proof}

\begin{proposition}[Composition]\label{prop:composition}
Let $f\colon S^{2n-1}\to S^n$, let $g\colon S^{2n-1}\to S^{2n-1}$ have degree $m$ and let $\varphi\colon S^n\to S^n$ have degree $k$. Then
\[
\Hopf(f\circ g)=m\,\Hopf(f),\qquad \Hopf(\varphi\circ f)=k^2\,\Hopf(f).
\]
\end{proposition}

\begin{proof}
Let $\hat g\colon D^{2n}\to D^{2n}$, $\hat g(tu)=t\,g(u)$, be the cone on $g$. It maps $S^{2n-1}$ to itself by $g$, and as a map of pairs $(D^{2n},S^{2n-1})\to(D^{2n},S^{2n-1})$ it has degree $m$, by the five lemma applied to the long exact sequences of the pairs and the boundary isomorphism $H^{2n}(D^{2n},S^{2n-1})\cong H^{2n-1}(S^{2n-1})$. The identity of $S^n$ together with $\hat g$ defines a map $\Phi\colon C_{f\circ g}\to C_f$, since a point $u\in\partial D^{2n}$ is attached at $f(g(u))$ in $C_{f\circ g}$ and $\hat g(u)=g(u)$ is attached at $f(g(u))$ in $C_f$:
\[
\begin{tikzcd}
S^{2n-1}\arrow[r,"g"]\arrow[d,hook] & S^{2n-1}\arrow[d,hook]\arrow[r,"f"] & S^n\\
D^{2n}\arrow[r,"\hat g"] & D^{2n} &
\end{tikzcd}
\qquad\leadsto\qquad
\Phi\colon C_{f\circ g}=S^n\cup_{fg}D^{2n}\longrightarrow S^n\cup_fD^{2n}=C_f .
\]
Then $\Phi^*\alpha_f=\alpha_{fg}$ and $\Phi^*\beta_f=m\beta_{fg}$, so applying $\Phi^*$ to $\alpha_f^2=\Hopf(f)\beta_f$ gives $\alpha_{fg}^2=m\Hopf(f)\beta_{fg}$. For the second identity let $\Psi\colon C_f\to C_{\varphi f}$ be $\varphi$ on $S^n$ and the identity on $D^{2n}$. Then $\Psi^*\alpha_{\varphi f}=k\alpha_f$ and $\Psi^*\beta_{\varphi f}=\beta_f$, and applying $\Psi^*$ to $\alpha_{\varphi f}^2=\Hopf(\varphi f)\beta_{\varphi f}$ gives $k^2\Hopf(f)\beta_f=\Hopf(\varphi f)\beta_f$.
\end{proof}

\begin{proposition}\label{prop:whiteheadH}
If $n$ is even, then $\Hopf([\iota_n,\iota_n])=\pm2$.
\end{proposition}

\begin{proof}
Write $S^n\times S^n=(S^n\vee S^n)\cup_wD^{2n}$ as in Section~\ref{ch:foundations}, and let $q\colon S^n\times S^n\to C=C_{\nabla\circ w}$ be the fold map $\nabla$ on $S^n\vee S^n$ and the identity on the top cell; it is well defined for the same reason as $\Phi$ above. Let $a,b\in H^n(S^n\times S^n)$ be the pull-backs of the generator along the two projections. Restriction to the $n$-skeleton is an isomorphism on $H^n$ for both spaces, and $\nabla^*$ sends the generator to $a+b$; hence $q^*\alpha=a+b$, where $\alpha=\alpha_{\nabla w}$. On the top cell $q$ is the identity, so $q^*\beta=\pm ab$, since $ab$ generates $H^{2n}(S^n\times S^n)$ by the K\"unneth theorem. Now $a^2=b^2=0$, as they are pulled back from $H^{2n}(S^n)=0$, and $ba=(-1)^{n^2}ab=ab$ since $n$ is even, so
\[
q^*(\alpha^2)=(a+b)^2=ab+ba=2ab=\pm2\,q^*\beta .
\]
As $q^*$ is an isomorphism on $H^{2n}$, $\alpha^2=\pm2\beta$.
\end{proof}

\subsection{Hopf invariant one}\label{sec:hone}

\begin{proof}[Proof of Theorem~\ref{thm:B}\textup{(i)}]
By Proposition~\ref{prop:cone}(ii) the mapping cone of $h_\K$ is homeomorphic to $\KP^2$, a closed simply connected $2d$-manifold, by a homeomorphism that is the identity on $S^d$ and of degree $\pm1$ on the top cell. Its cohomology is free, and by Poincar\'e duality (input (I2), in the form \cite[Proposition~3.38]{Hat02}) the cup product pairing $H^d(\KP^2)\times H^d(\KP^2)\to H^{2d}(\KP^2)\cong\Z$ is nonsingular. Since $H^d$ is generated by $\alpha$, this says $\alpha^2=\pm\beta$.
\end{proof}

\begin{remark}
The proof uses that $C_{h}$ is a manifold, and so it applies to any fibre bundle $p\colon S^{2d-1}\to S^d$ with fibre $S^{d-1}$, $d\ge2$: coning the fibres turns $p$ into a bundle with fibre $D^d$, which is a manifold with boundary $S^{2d-1}$, and gluing $D^{2d}$ gives a closed manifold homeomorphic to $C_p$. With Adams' theorem (I4) this gives the following statement, quoted in the Introduction.
\end{remark}

\begin{corollary}\label{cor:spherebundles}
If $S^{d-1}\to S^{2d-1}\xrightarrow{\,p\,}S^d$ is a fibre bundle with $d\ge2$, then $\Hopf(p)=\pm1$ and $d\in\{2,4,8\}$.
\end{corollary}

\subsection{Power maps}\label{sec:powers}

By Artin's theorem the powers $x^m$ of an octonion are well defined, since $x$ generates an associative subalgebra. For $d\in\{2,4,8\}$ let
\[
p_m\colon S^{d-1}\to S^{d-1},\qquad p_m(x)=x^m ,
\]
which takes values in the unit sphere by \ref{D:norm}. Every $x\in S^{d-1}\setminus\{\pm1\}$ can be written uniquely as $x=\cos\theta+u\sin\theta$ with $0<\theta<\pi$ and $u$ a unit vector orthogonal to $1$; such a $u$ satisfies $u^2=-1$ by \ref{D:conj}, so $\C_u=\R\oplus\R u$ is a subalgebra isomorphic to $\C$, and $x^m=\cos m\theta+u\sin m\theta$.

\begin{lemma}\label{lem:powerderiv}
Let $x=\cos\theta+u\sin\theta$ as above and let $W_u\subset\K$ be the orthogonal complement of $\C_u$. Then $W_u$ is tangent to $S^{d-1}$ at $x$ and at $x^m$, and
\[
Dp_m(x)(w)=\frac{\sin m\theta}{\sin\theta}\,w\qquad(w\in W_u),\qquad Dp_m(x)(ux)=m\,u\,x^m .
\]
\end{lemma}

\begin{proof}
The second formula is the derivative of $e^{i\theta}\mapsto e^{im\theta}$ in $\C_u$, since $ux=\frac{d}{d\theta}x$. For the first, $x$ and $w$ lie in a subalgebra that is associative by Artin's theorem, so $Dp_m(x)(w)=\sum_{j=0}^{m-1}x^jwx^{m-1-j}$, computed by expanding $(x+tw)^m$. Since $w\perp1$ and $w\perp u$, \ref{D:inner} gives $\Re(w)=0$ and $\Re(u\bar w)=0$, hence $\bar w=-w$ and $uw=-wu$ (the latter because $uw+wu=-(u\bar w+w\bar u)=-2\langle u,w\rangle=0$). Therefore $wx=\bar xw$, and $x^jwx^{m-1-j}=x^j\bar x^{m-1-j}w=x^{2j-m+1}w$ because $\bar x=x^{-1}$ in $\C_u$. Finally
\[
\sum_{j=0}^{m-1}x^{2j-m+1}=\sum_{j=0}^{m-1}\bigl(\cos(2j-m+1)\theta+u\sin(2j-m+1)\theta\bigr)=\frac{\sin m\theta}{\sin\theta},
\]
the sine terms cancelling in pairs $j\leftrightarrow m-1-j$. The vectors $w\in W_u$ are orthogonal to $x$ and to $x^m$, both of which lie in $\C_u$, so they are tangent to the sphere at both points.
\end{proof}

\begin{proposition}\label{prop:powerdegree}
For $d\in\{2,4,8\}$ and $m\ge1$ the power map $p_m$ has degree $m$.
\end{proposition}

\begin{proof}
Let $y=\cos\varphi+v\sin\varphi$ with $0<\varphi<\pi$, a non-real point of $S^{d-1}$. If $x^m=y$, then $y\in\C_x$, and $\C_x=\C_y=\R\oplus\R y$ because a two-dimensional subalgebra containing the non-real element $y$ is spanned by $1$ and $y$. So the solutions of $x^m=y$ are the $m$ roots of $y$ in the circle $\C_y\cap S^{d-1}$, and none of them is real. Write such a root as $x=\cos\theta+v\sin\theta$ with $\theta\in(-\pi,\pi)\setminus\{0\}$; the formulas of Lemma~\ref{lem:powerderiv} are unchanged under $(u,\theta)\mapsto(-u,-\theta)$, so they hold with $u=v$. We have $\sin\theta\ne0$ and $\sin m\theta=\pm\sin\varphi\ne0$. By Lemma~\ref{lem:powerderiv}, in the bases $(vx,w_1,\dots,w_{d-2})$ of $T_xS^{d-1}$ and $(vy,w_1,\dots,w_{d-2})$ of $T_yS^{d-1}$, with $(w_j)$ a fixed basis of $W_v$, the derivative is $\operatorname{diag}(m,c,\dots,c)$ with $c=\sin m\theta/\sin\theta\ne0$. Its determinant $m\,c^{d-2}$ is positive, since $d-2\in\{0,2,6\}$ is even. Both bases are positively oriented or both negatively: the frames $(x,vx,w_1,\dots)$ and $(y,vy,w_1,\dots)$ of $\R^d$ differ by a rotation of the plane $\C_v$ that fixes $W_v$. So $y$ is a regular value with $m$ preimages, each counted with sign $+1$.
\end{proof}

\begin{remark}
The scalar $c=\sin m\theta/\sin\theta$ vanishes at $\theta=k\pi/m$, $0<k<m$, where $x^m=\pm1$; the critical points of $p_m$ for $d\ge4$ form the $(d-2)$-spheres $\{\cos(k\pi/m)+u\sin(k\pi/m):u\perp1,\ |u|=1\}$, which $p_m$ collapses to the points $(-1)^k$. For $d=2$ there are no critical points.
\end{remark}

\begin{proposition}\label{prop:gm}
For $d\in\{2,4,8\}$ and $m\ge1$ the map $g_m\colon S^{2d-1}\to S^{2d-1}$, $g_m(a,b)=(a^m,b)/|(a^m,b)|$, has degree $m$, and it fixes the point $(1,0)$. The map $\kappa(a,b)=(\bar a,b)$ has degree $-1$.
\end{proposition}

\begin{proof}
Write $a=\rho x$ with $\rho=|a|$ and $|x|=1$. For $t\in[0,1]$ the maps
\[
F_t(\rho x,b)=\bigl(\rho^{(1-t)m+t}x^m,\;b\bigr)
\]
are continuous on $\K^2$ (the exponent is at least $1$) and vanish only at $0$, so $F_t/|F_t|$ is a homotopy of self-maps of $S^{2d-1}$ fixing $(1,0)$, from $g_m$ to $g'(a,b)=(|a|\,p_m(a/|a|),\,b)$, extended by $g'(0,b)=(0,b)$; note that $(|a|\,p_m(a/|a|),b)$ already has norm one. On the open dense subset $\{a\neq0,\ b\neq0\}$ of $S^{2d-1}$ we use the diffeomorphism
\[
(0,\tfrac{\pi}{2})\times S^{d-1}\times S^{d-1}\to\{a\ne0,b\ne0\},\qquad (s,x,z)\mapsto(x\cos s,\;z\sin s),
\]
in which $g'$ becomes $(s,x,z)\mapsto(s,p_m(x),z)$. A regular value $y$ of $p_m$ gives a regular value $(y\cos s,z\sin s)$ of $g'$, whose preimages $(x\cos s,z\sin s)$, $x^m=y$, are $m$ points at which the Jacobian in these coordinates has the sign of $\det Dp_m(x)$, that is $+1$. The same chart is used in source and target, so the sign relative to the orientation of $S^{2d-1}$ is also $+1$, and $\deg g_m=\deg g'=m$. The map $\kappa$ is the restriction of a linear map of $\R^{2d}$ that negates the $d-1$ imaginary coordinates of $a$, so its degree is $(-1)^{d-1}=-1$.
\end{proof}

\begin{proof}[Proof of Theorem~\ref{thm:B}\textup{(ii)} and \textup{(iii)}]
Part (ii) is Propositions~\ref{prop:powerdegree} and \ref{prop:gm}. Since $g_m$ is a based map of degree $m$, Lemma~\ref{lem:precompose} gives $[h_\K\circ g_m]=m[h_\K]$ in $\pi_{2d-1}(S^d)$, and Proposition~\ref{prop:composition} gives $\Hopf(h_\K\circ g_m)=m\Hopf(h_\K)=\pm m$ and $\Hopf(h_\K\circ g_m\circ\kappa)=\mp m$. Because $h_\K(\lambda v)=\lambda^2h_\K(v)$ for real $\lambda$,
\begin{equation}\label{eq:hgm}
h_\K\circ g_m(a,b)=\frac{\bigl(2a^m\bar b,\;|a|^{2m}-|b|^2\bigr)}{|a|^{2m}+|b|^2},
\end{equation}
whose numerator and denominator are polynomials in the real coordinates of $(a,b)$. With the constant map for $0$, every integer is the Hopf invariant of such a map.
\end{proof}

\begin{corollary}\label{cor:image}
The image of $\Hopf\colon\pi_{2n-1}(S^n)\to\Z$ is $0$ for $n$ odd, $\Z$ for $n=2,4,8$, and $2\Z$ for every other even $n$.
\end{corollary}

\begin{proof}
The first case is Proposition~\ref{prop:hodd}. For $n$ even the image is a subgroup containing $2$ by Proposition~\ref{prop:whiteheadH}, so it is $\Z$ or $2\Z$; it is $\Z$ for $n=2,4,8$ by Theorem~\ref{thm:B}(i), and $2\Z$ otherwise by Adams' theorem (I4).
\end{proof}

\begin{remark}\label{rem:whitehead-hopf}
For $n=2$ the group $\pi_3(S^2)\cong\Z$ (Theorem~\ref{thm:C}) is detected by $\Hopf$, so Proposition~\ref{prop:whiteheadH} gives $[\iota_2,\iota_2]=\pm2[h]$. For $n=4$ it gives $[\iota_4,\iota_4]\equiv\pm2[h_\Hq]$ modulo the torsion subgroup $\Z/12$ of $\pi_7(S^4)$; the exact formula $[\iota_4,\iota_4]=2\nu_4-E\nu'$ in Toda's notation is \cite[Chapter~V]{Tod62}. Proposition~\ref{prop:composition} also shows that $\Hopf(h\circ\varphi)=k$ for a self-map $\varphi$ of $S^3$ of degree $k$, while $\Hopf(\varphi'\circ h)=k^2$ for a self-map $\varphi'$ of $S^2$ of degree $k$; both appear again in Section~\ref{ch:linking}, as $[z^k:w]$ and $[z^k:w^k]$.
\end{remark}
